% Independent CI run 38363154062; exact rational polynomial root-stage scope.
\paragraph{Strong exact root-isolation comparator.}
The incremental event update previously held the exact SymPy
\texttt{real\_roots} engine fixed, comparing only full per-cell
re-evaluation with incident-pair updates. We separately compared that
root engine with SymPy's established exact rational isolating
intervals at width at most $10^{-8}$, using identical constructed
rational pair polynomials in 3 and 6 spectral dimensions. All
intervals contained the corresponding exact algebraic roots with
their multiplicities; an exact root-count check also covered the
interior $\lambda\in(0,1)$. The intervals algorithm is a strong
\emph{existing} comparator, not our contribution.

\begin{table}[h!]\centering\small
\caption{Exact root-stage timing on the public CI machine, median
seconds over four alternating-order repeats. Each method isolates all
real roots of every polynomial, not an end-to-end AUC path.
Times depend on the machine and fixed constructed examples.}
\label{tab:p4-root-isolator}
\begin{tabular}{llrrr}
\toprule
Dimension & Pattern & Pairs & Rational intervals & \texttt{real\_roots}\\
\midrule
3 & Crossing & 16 & .00416 & .01081\\
6 & Crossing & 16 & .00888 & .03562\\
6 & No event & 16 & .04569 & .02819\\
\bottomrule
\end{tabular}
\end{table}

The rational isolator is already faster at the \emph{root stage} in
the crossing examples, while it is slower in the six-dimensional
no-event control. Thus the previous speedup over full AUC
re-evaluation cannot be interpreted as superiority to a strong
root-isolation method. A complete strict algorithm must combine
root isolation, equality and ordering of roots across pairs, tie
scores, incident updates and floating input error control in one
fair end-to-end benchmark. The comparator code and full eight-case
receipt ran in
\href{https://github.com/lindicaphxag-tech/kaggle/actions/runs/38363154062}{the successful independent CI workflow}.

\paragraph{Complete event map with the same strong isolator.}
We next used these exact rational isolating intervals in \emph{both}
complete event-map algorithms. A rational interval that lay wholly
outside $(0,1)$ was rejected without computing its algebraic root;
for any nonrational interval overlapping the feasible parameter range,
the solver falls back to \texttt{real\_roots}. We compared full
per-event pair re-evaluation against incremental updates of incident
pairs. The two outputs agreed exactly with one another and with the
original exact event-map implementation, including endpoint ties,
zero polynomials, coincident roots, even multiplicities and an
irrational-root check.

On the independent CI machine, the six-dimensional 16-pair,
32-event constructed example took .14495~s for full re-evaluation
and .02065~s for the incremental method (approximately 7.0 times
faster), using the \emph{same} stronger root engine. The
three-dimensional 32-pair, 64-event case took .08380 versus
.01856~s (approximately 4.5 times faster). In the six-dimensional
16-pair \emph{no-event} control the times were .04998 and
.05000~s, so there is no meaningful gain. This isolates a
conditional event-update advantage on these constructed rational
inputs; it does not imply faster roots, a universal speedup,
floating-input correctness, or a real-EEG performance effect. The
complete six-case workflow is
\href{https://github.com/lindicaphxag-tech/kaggle/actions/runs/38364906164}{available with an exact parity check}.
